%%%%%%%%%%%%%%%%%%%%%%%%%%%%%%%%%%%%%%%%%%%%%%%%%%%%%%%%%%%%
\section{Tabellen}%
\label{sec:Tabellen}
%%%%%%%%%%%%%%%%%%%%%%%%%%%%%%%%%%%%%%%%%%%%%%%%%%%%%%%%%%%%
Für eine ausführliche Erläuterung auch über gute und schlechte Tabellen
und deren Gestaltung empfiehlt sich die Lektüre des entsprechenden Artikels im \LaTeX{}"=Kompendium auf Wikibooks%
\footnote{\url{https://de.wikibooks.org/wiki/LaTeX-Kompendium:_Tabellen}}
sowie die Dokumentation des \pkg{booktabs}-Pakets \cite{Fear2005}.
\Ua stellt dieses Paket die Befehle
\begin{itemize*}
	\item \lstinline|\toprule|
	\item \lstinline|\midrule|
	\item \lstinline|\bottomrule|
\end{itemize*}
zur Verfügung.

Der Code für eine einfache Tabelle ist im \cref{lst:tabellenbeispiel} angeführt.
Damit alle Tabellen einheitlich aussehen, empfiehlt es sich,
überall das gleiche Makro zur Parameteranpassung zu verwenden.
Dazu könnte beispielsweise das vordefinierte Makro
\lstinline|\tablestylecommon| verwendet werden, welches in der Datei
\printfilepathquoted{./preambel/TableCommands.tex} definiert ist.
Zu beachten ist, dass die Anpassungen an den Tabellenparametern wie \zB
\begin{itemize*}
\item die Verkleinerung der Schriftgröße,
\item die Entfernung der seitlichen Ränder sowie
\item die Erhöhung der Abstandes zwischen den einzelnen Tabellenzeilen
\end{itemize*}
innerhalb einer Umgebung gesetzt werden müssen, damit sie lokal wirken und
sich nicht auf die nachfolgenden Tabellen erstrecken.
Wird die Tabelle nicht innerhalb einer \printkeyword{table}-Umgebung
\oae gesetzt, muss man eine eigene Umgebung mit
\lstinline|\begingroup| ... \lstinline|\endgroup|
definieren.

\begin{latex}[caption={Einfache Tabelle in \LaTeX},label={lst:tabellenbeispiel}]
\begin{table}%
	\tablestylecommon% Aussehen anpassen
	\caption{Dateien der Vorlage}%
	\label{tab:tabellenbeispiel}%
	\begin{tabular}{l l l}%
		\toprule%
		Datei       &  Bedeutung    &  Benutzerinteraktion \\%
		\midrule%
		Diss.tex  &  Hauptdatei   &  nein     \\%
		images/   &  Bilder       &  ja       \\%
		content/  &  Kapitel      &  ja       \\%
		\bottomrule%
	\end{tabular}%
\end{table}
\end{latex}

Das Ergebnis sieht man in \cref{tab:tabellenbeispiel}.

\begin{table}%
	\tablestylecommon% Aussehen anpassen
	\caption{Dateien der Vorlage}%
	\label{tab:tabellenbeispiel}%
	\begin{tabular}{l l l}%
		\toprule%
		Datei       &  Bedeutung    &  Benutzerinteraktion \\%
		\midrule%
		Diss.tex  &  Hauptdatei   &  nein     \\%
		images/   &  Bilder       &  ja       \\%
		content/  &  Kapitel      &  ja       \\%
		\bottomrule%
	\end{tabular}%
\end{table}

Ist die Tabelle seitenbreit, sollte statt der \printkeyword{tabular}-Umgebung die \printkeyword{tabularx}-Umgebung verwendet werden (\cref{tab:AlternierendeZeilenfarben}).
Mit dem \printkeyword{tablealtcolored}-Makro kann man die einzelnen Zeilen mit alternierender Farbe einfärben.
Die Anpassungsmöglichkeiten für das Aussehen der Tabellen befinden sich in der Datei \printfilepathquoted{./preambel/TableCommands.tex}.

%% Breite der ersten und zweiten Spalte auf Inhalt anpassen --> "l".
%% Breite der dritten Spalte automatisch bestimmen --> "X" (oder "x", "F", "y", "Y")
%% "X" = Blocksatz (mit Silbentrennung),
%% "x" = linksbündiger Flattersatz ohne Silbentrennung
%% "F" = linksbündiger Flattersatz mit Silbentrennung (KSP-Verlag: unerwünscht),
%% "y" = rechtsbündiger Flattersatz ohne Silbertrennung
%% "Y" = rechtsbündiger Flattersatz mit Silbertrennung
\begin{table}%
\tablestylealtcolored%
\caption{Seitenbreite Tabelle mit alternierender Zeilenfarbe}%
\label{tab:AlternierendeZeilenfarben}%
	\begin{tabularx}{\linewidth}{llX}%
		\toprule%
		\rowcolor{tableheadcolor}%
		\tableheadfont Tabellenkopf &	\tableheadfont Tabellenkopf & \tableheadfont Tabellenkopf 
		\tabularnewline%
		\midrule%
		Inhalt  & Etwas längerer Inhalt & Weiterer Inhalt \tabularnewline
		Inhalt  & Etwas längerer Inhalt & Weiterer Inhalt \tabularnewline
		Inhalt  & Etwas längerer Inhalt & Weiterer Inhalt \tabularnewline
		%% Zwischenkopf ---------------------------------------------
		\rowcolor{tablesubheadcolor}%
		\multicolumn{3}{l}{\bfseries\color{KITblue} Zwischenkopfzeile}%
		\tabularnewline%
		%%-----------------------------------------------------------
		Inhalt  & Etwas längerer Inhalt & Weiterer Inhalt \tabularnewline
		Inhalt  & Etwas längerer Inhalt & Weiterer Inhalt \tabularnewline
		Inhalt  & Etwas längerer Inhalt & Weiterer Inhalt \tabularnewline
		%% Zwischenkopf ---------------------------------------------
		\rowcolor{tablesubheadcolor}%
		\multicolumn{3}{l}{\bfseries\color{KITgreen} Zwischenkopfzeile}%
		\tabularnewline
		%%-----------------------------------------------------------
		Inhalt  & Etwas längerer Inhalt & Weiterer Inhalt \tabularnewline%
		Inhalt  & Etwas längerer Inhalt & Weiterer Inhalt \tabularnewline%
		\bottomrule%
	\end{tabularx}%
\end{table}

%% ------------------------------------------------------------
Lange Tabellen, die über mehrere Seiten gehen und umbrochen werden sollen, können mit
\lstinline|\LTXtable{\textwidth}{Datei}|
eingebunden werden, wobei die Tabelle in eine eigene Datei ausgelagert werden muss.
Ein Beispiel dafür sieht man in \cref{tab:MehrseitigeTabelle}.

% Angaben zur Tabellenstil innerhalb einer eigenen Umgebung, damit diese nur lokal wirken
\begingroup%
\tablestylealtcolored%
\setlength{\topskip}{0pt}%
\setlength{\LTpre}{0pt}%
\setlength{\LTpost}{0pt}%
\LTXtable{\textwidth}{tables/LongTableExample.tex}
\endgroup%

Sollte eine Tabelle so breit sein, dass sie nicht mehr horizontal auf
eine Seite passt, so ist es natürlich möglich, diese mithilfe des Pakets
\printkeyword{rotfloat} \parencite{Sommerfeldt2004} in eine
\printkeyword{sidewaystable} statt in eine \printkeyword{table}-Umgebung zu setzen.
Also so:
\begin{latex}[caption={Gedrehte Tabelle},label={lst:rotated-table}]
\begin{sidewaystable}
  \raggedleft% rechtsbündig, um die Tabelle oben auszurichten
  \begin{tabular}{...}%
    ...
  \end{tabular}%
  \caption{Bezeichnung}%
  \label{Referenzmarke}%
\end{sidewaystable}%
\end{latex}

Das Ergebnis sieht man in \cref{tab:ex-sideways}.

%% Um 90° gedrehte Tabelle
%
\begin{sidewaystable}[p]
\raggedleft%% Rechtsbündig, damit die Tabelle oben ausgerichtet wird (Forderung KSP)
\renewcommand{\arraystretch}{1.3}%% Etwas mehr Abstand zwischen den Tabellenzeilen
%% Tabelle mit einer festen Breite (hier: so breit wie die Seitenhöhe).
%% Spaltendefinition:
%% Erste Spalte vom Typ "l": Spalte so breit, wie der Inhalt, linksbündige Ausrichtung
%% Zweite Spalte vom Typ "L": Spalte mit einer festen Breitenangabe, linksbündiger Flattersatz. Inhalt wird umbrochen, wenn zu lang
%% 2 Spalten vom Typ "x":  die Breite wird automatisch berechnet
%% in Abhängigkeit von der Gesamtbreite der Tabelle und Breite anderer Spalten,
%% Inhalt wird umbrochen, wenn zu lang
%% Format: linksbündiger Flattersatz ohne Silbentrennung, 
%% Weitere Möglichkeiten:
%% "X" = Blocksatz (mit Silbentrennung)
%% "F" = linksbündiger Flattersatz mit Silbentrennung (KSP-Verlag: unerwünscht),
%% "y" = rechtsbündiger Flattersatz ohne Silbertrennung
%% "Y" = rechtsbündiger Flattersatz mit Silbertrennung
\begin{tabularx}{\linewidth}{l L{10em} *{2}x}%
	\toprule%
	\rowcolor{tableheadcolor}%
	\tableheadfont Spalte vom Typ \enquote{l} &	
	\tableheadfont Spalte vom Typ \enquote{L} & 
	\tableheadfont Spalte vom Typ \enquote{x} & 
	\tableheadfont Spalte vom Typ \enquote{x}
	\tabularnewline%
	\midrule%
		Spaltenbreite = Breite des Inhalts & 
		Spalte mit der festen Breite von 10\,em & 
		Spaltenbreite wird automatisch bestimmt. Inhalt wird bei Bedarf umbrochen. & 
		Spaltenbreite wird automatisch bestimmt. Inhalt wird bei Bedarf umbrochen.
	\tabularnewline
	%%-----------------------------------------------------------
	Nächste Zeile & Etwas längerer Inhalt & Weiterer Inhalt & und noch etwas \tabularnewline
	Weitere Zeile & Etwas längerer Inhalt & Weiterer Inhalt & und noch etwas \tabularnewline
	%% Zwischenkopf ---------------------------------------------
	\rowcolor{tablesubheadcolor}%
	\multicolumn{4}{l}{\bfseries\color{KITblue} Zwischenkopfzeile, welche über mehrere Spalten geht}%
	\tabularnewline%
	%%-----------------------------------------------------------
	Inhalt  & Etwas längerer Inhalt & Weiterer Inhalt & und noch etwas \tabularnewline
	Inhalt  & Etwas längerer Inhalt & Weiterer Inhalt & und noch etwas \tabularnewline
	Inhalt  & Etwas längerer Inhalt & Weiterer Inhalt & und noch etwas \tabularnewline
	%% Zwischenkopf ---------------------------------------------
	\rowcolor{tablesubheadcolor}%
	\multicolumn{4}{l}{\bfseries\color{KITgreen} Zwischenkopfzeile, welche über mehrere Spalten geht}%
	\tabularnewline
	%%-----------------------------------------------------------
	Inhalt  & Etwas längerer Inhalt & Weiterer Inhalt & und noch etwas \tabularnewline
	Inhalt  & Etwas längerer Inhalt & Weiterer Inhalt & und noch etwas \tabularnewline
	\bottomrule%
\end{tabularx}%
\caption{Beispiel für eine breite, gedrehte Tabelle}%
\label{tab:ex-sideways}%
\end{sidewaystable}